\documentclass[10pt,a4paper]{article}

\usepackage[margin=20mm]{geometry}
\usepackage{amsmath,amssymb,amsthm}

\theoremstyle{definition}
\newtheorem*{definition*}{Definition}
\newtheorem*{problem*}{Problem}

\setlength{\parindent}{0pt}
\setlength{\parskip}{5pt}

\begin{document}

\begin{center}
  {\Large\bfseries Oracle Comparator Circuits versus Oracle Catalytic Logspace}
\end{center}

\section{Definitions}

\begin{definition*}[Comparator circuits and $\mathsf{CC}$]
A comparator gate is the two-input, two-output Boolean function
\begin{equation}
  C:\{0,1\}^2\longrightarrow\{0,1\}^2,
  \qquad
  C(p,q)=(p\wedge q,p\vee q).
\end{equation}
A comparator circuit on $m$ wires is specified by a finite sequence
\begin{equation}
  (i_1,j_1),\ldots,(i_t,j_t),
\end{equation}
where gate $(i_k,j_k)$ applies $C$ to the current contents of wires $i_k$ and
$j_k$. Each gate replaces the two values by their minimum and maximum and does
not copy either input. Thus the model has fan-out one.

An annotated comparator circuit designates one output wire and labels each
input wire by a constant $0$ or $1$, an input variable $x_i$, or a negated
input variable $\neg x_i$. A family $\{C_n\}_{n\geq1}$ is
$\mathsf{AC}^0$-uniform when the description of $C_n$ is generated by a
function in uniform $\mathsf{FAC}^0$. The class $\mathsf{CC}$ consists of the
relations computed by $\mathsf{AC}^0$-uniform families of polynomial-size
annotated comparator circuits.
\end{definition*}

\begin{definition*}[Catalytic logspace and $\mathsf{CL}$]
A catalytic Turing machine is a deterministic Turing machine equipped with a
read-only input tape, an ordinary read-write work tape, and an additional
read-write auxiliary tape. For an input $x$ and an arbitrary string $a$
initially written on the auxiliary tape, let $M(x,a)$ denote the resulting
computation. The machine decides a language $L$ if, for every $x$ and $a$,
\begin{enumerate}
  \item $M(x,a)$ halts with the auxiliary tape restored exactly to $a$; and
  \item $M(x,a)$ accepts if and only if $x\in L$.
\end{enumerate}
For functions $S,S_a:\mathbb N\to\mathbb N$, define
\begin{equation}
  \mathsf{CSPACE}(S(n),S_a(n))
\end{equation}
as the class of languages decided by catalytic machines using $O(S(n))$ cells
of ordinary work tape and $O(S_a(n))$ cells of auxiliary tape on inputs of
length $n$. Write
\begin{equation}
  \mathsf{CSPACE}(S(n))
  :=
  \mathsf{CSPACE}\bigl(S(n),2^{O(S(n))}\bigr),
  \qquad
  \mathsf{CL}:=\mathsf{CSPACE}(\log n).
\end{equation}
\end{definition*}

\begin{definition*}[Native functional oracle for comparator circuits]
A comparator-circuit oracle is a length-preserving function
\begin{equation}
  \alpha:\{0,1\}^*\longrightarrow\{0,1\}^*,
  \qquad
  |\alpha(y)|=|y|.
\end{equation}
For each $n$, let $\alpha_n$ be the restriction of $\alpha$ to $\{0,1\}^n$.
An $\alpha_n$-gate in a comparator circuit selects $n$ distinct wires carrying
$y\in\{0,1\}^n$ and sends the $n$ bits of $\alpha_n(y)$ to $n$ distinct
outgoing wires. Each output is connected to at most one later gate input, so
the limited-fan-out condition is preserved.

For an input $X$ of length $n$, a relation $R(X,\alpha)$ belongs to
$\mathsf{CC}(\alpha)$ when it is computed by a polynomial-size uniform family
of annotated comparator circuits whose allowed gates are comparator gates,
negation gates, and $\alpha_n$-gates.
\end{definition*}

\begin{definition*}[Weak oracle for catalytic logspace]
Let $O:\{0,1\}^*\to\{0,1\}$ be a Boolean oracle. A $\mathsf{CL}^O$ machine has
a one-bit read-only oracle-output register, initially $0$, but no query tape.
An oracle call designates a subinterval $T$ of data already present on the
input or catalytic tape. If $y$ is the current content of $T$, the call performs
\begin{equation}
  \text{oracle-output}\leftarrow O(y)
\end{equation}
without modifying $T$. The two endpoints of $T$ are encoded using $O(\log n)$
bits of the machine's existing storage. For every input and every initial
catalyst, the machine must return the correct answer and restore the catalyst
exactly.
\end{definition*}
\begin{definition*}[Exact convention for $\mathsf{CC}^{O}$ with a Boolean
oracle]

To compare comparator circuits with $\mathsf{CL}^{O}$ relative to one and the
same Boolean oracle, fix
\begin{equation}
                 O:\{0,1\}^{*}\longrightarrow\{0,1\}
\end{equation}
and associate with it the length-preserving function $\widehat O$ given by
\begin{equation}
  \widehat O(\varepsilon)=\varepsilon,
  \qquad
  \widehat O(by)=\bigl(b\oplus O(y)\bigr)y
  \quad
  (b\in\{0,1\},\;y\in\{0,1\}^{*}).
\end{equation}
Thus $|\widehat O(z)|=|z|$. On $r=m+1$ wires, a
$\widehat O_r$-gate acts as
\begin{equation}
  (b,y_1,\ldots,y_m)
  \longmapsto
  \bigl(b\oplus O(y_1\cdots y_m),y_1,\ldots,y_m\bigr).
\end{equation}
It exposes $O(y)$ when its answer wire is initialized to $b=0$, preserves the
query bits, has equally many inputs and outputs, and creates no fan-out.

For this paper, define
\begin{equation}
                         \boxed{\mathsf{CC}^{O}:=\mathsf{CC}(\widehat O)}.
\end{equation}
Equivalently, a language $A$ belongs to $\mathsf{CC}^{O}$ if there is an
$\mathsf{AC}^{0}$-uniform polynomial-size family of annotated comparator
circuits, independent of the values of $O$, using comparator gates, negation
gates, and the corresponding $\widehat O_n$-gates, whose designated output is
$1$ exactly on the strings in $L$. This reversible Boolean embedding is the
comparison convention adopted here.
\end{definition*}



\section{Best Known Result}

The unrelativized containments are
\begin{equation}
  \mathsf{NL}\subseteq\mathsf{CC}\subseteq\mathsf P,
  \qquad
  \mathsf{NL}\subseteq\mathsf{CL}\subseteq\mathsf{ZPP}.
\end{equation}

In the native functional-oracle model for comparator circuits, there are
suitable oracle settings $\alpha$ and $\beta$, not necessarily the same, for
which
\begin{equation}
  \mathsf{CC}(\alpha)\not\subseteq\mathsf{NC}(\alpha),
  \qquad
  \mathsf{NC}^3(\beta)\not\subseteq\mathsf{CC}(\beta).
\end{equation}
These separations remain valid when the functional oracles are restricted to
be strictly $1$-Lipschitz. The relativization of $\mathsf{NC}$ used in this
model controls the nesting depth of oracle gates and preserves containments
such as
\begin{equation}
  \mathsf{NC}^1(\alpha)\subseteq\mathsf L(\alpha).
\end{equation}
These results do not separate the common Boolean-oracle classes appearing in
the problem below, because an arbitrary length-preserving functional oracle
need not arise from a Boolean oracle.

For oracle catalytic logspace, there are oracle worlds in which the class
collapses to ordinary logspace and others in which, for a suitable $O$,
\begin{equation}
  \mathsf{CL}^O=\mathsf{PSPACE}^O.
\end{equation}
The strongest result stated in the source is
\begin{equation}
  \exists O\qquad \mathsf{CL}^O=\mathsf{EXP}^O,
\end{equation}
even under the weak no-query-tape convention above.

\section{Problem Formalization}

Prove that there are oracles $O_1$ and $O_2$ such that
\begin{equation}
  \mathsf{CC}^{O_1}\not\subseteq\mathsf{CL}^{O_1}
  \qquad\text{and}\qquad
  \mathsf{CL}^{O_2}\not\subseteq\mathsf{CC}^{O_2}.
\end{equation}

Then prove the stronger version: there exists a single oracle $O$ such that
\begin{equation}
  \mathsf{CC}^{O}\not\subseteq\mathsf{CL}^{O}
  \qquad\text{and}\qquad
  \mathsf{CL}^{O}\not\subseteq\mathsf{CC}^{O}.
\end{equation}


\end{document}